% Fragment -- inclus de main.tex prin \input
\clearpage
\setpartlabel{themeHtml}{\faCode}{Faza 1: HTML}
\setcounter{section}{0}
\phantomsection
\addcontentsline{toc}{part}{\protect\textbf{Faza 1 \textendash{} HTML}}
\handout{Faza 1 \textendash{} HTML}{Structura documentului. Tag-uri semantice. Accesibilitate.}

HTML (\emph{HyperText Markup Language}) este limbajul de marcare folosit pentru descrierea \textbf{structurii} unui document web. Browserul citeste documentul de sus in jos, construieste in memorie un arbore de noduri (DOM-ul), si afiseaza paginea pe ecran.

In arhitectura unei pagini web, responsabilitatile sunt impartite intre trei tehnologii:

\begin{itemize}[itemsep=0pt]
  \item \textbf{HTML} \textendash{} ce contine pagina (titluri, paragrafe, imagini, formulare).
  \item \textbf{CSS} \textendash{} cum arata pagina (culori, dimensiuni, spatiere, layout).
  \item \textbf{JavaScript} \textendash{} cum se comporta pagina (raspuns la click, animatii, validare).
\end{itemize}

Capitolul defineste anatomia minimala, tag-urile semantice principale, atributele obligatorii pentru accesibilitate, si patternurile uzuale pentru linkuri si formulare.

\bigskip

\begin{outcomes}
Dupa capitolul asta vei \textcommabelow{s}ti:
\begin{itemize}[itemsep=1pt]
  \item Scrie boilerplate-ul HTML5 corect, cu cele 3 meta tags obligatorii.
  \item Distinge tag-urile semantice de \code{<div>} \textcommabelow{s}i alege tag-ul potrivit per zona.
  \item Construie\textcommabelow{s}te formulare accesibile cu \code{<label for>} \textcommabelow{s}i \code{<input id>}.
  \item Aplica atributele \code{alt}, \code{lang}, \code{href} cu valori corecte conform WCAG.
  \item Recunoa\textcommabelow{s}te pattern-urile pentru linkuri externe sigure (\code{rel="noopener"}).
  \item \faRocket\ \emph{Bonus self-study:} \code{<picture>} cu fallback de format, \code{loading="lazy"} + \code{decoding="async"}, ARIA (\code{aria-label}, \code{aria-current}).
\end{itemize}
\end{outcomes}

\begin{whymatters}
HTML este temelia oricarei pagini web. Browserele, cititoarele de ecran \textcommabelow{s}i motoarele de cautare interpreteaza pagina pornind de la structura HTML. O structura semantica gre\textcommabelow{s}ita conduce la trei consecin\textcommabelow{t}e directe: SEO penalizat, persoanele cu dizabilita\textcommabelow{t}i nu pot folosi pagina, codul devine ilizibil dupa cateva luni. Toate fazele urmatoare (CSS, JS, deploy, audit) presupun ca structura HTML este corecta.
\end{whymatters}

% =========================================================
\section{Anatomia documentului}

Un document HTML5 contine cinci elemente structurale obligatorii. Pe linia 1, declaratia \code{<!DOCTYPE html>} indica browserului ca documentul respecta specificatia HTML5; in absenta sa, browserul intra in modul ,,quirks'' care reproduce comportamentul browserelor din anii 2000.

\begin{lstlisting}[style=html, caption={Boilerplate HTML5 conform specificatiei W3C}]
<!DOCTYPE html>
<html lang="ro">
<head>
  <meta charset="UTF-8">
  <meta name="viewport" content="width=device-width, initial-scale=1.0">
  <title>Numele Prenumele -- Student CS, Iasi</title>
</head>
<body>
  <!-- continut vizibil -->
</body>
</html>
\end{lstlisting}

\begin{itemize}
  \item \code{<html lang="ro">} \textendash{} elementul radacina. Atributul \code{lang} influenteaza pronuntia in cititoare de ecran, ranking-ul in cautari localizate, si detectia limbii sursa pentru traducere automata.
  \item \code{<head>} \textendash{} contine metadatele paginii. Continutul nu apare pe ecran, dar este consumat de browser, motoare de cautare si platforme sociale.
  \item \code{<body>} \textendash{} contine elementele vizibile pe ecran. Tot ce vede utilizatorul se afla in interiorul acestui tag.
\end{itemize}

% =========================================================
\section{Sintaxa tag-urilor}

Un element HTML are una sau mai multe parti: tag de deschidere, continut, tag de inchidere. Atributele se declara in tag-ul de deschidere, in formatul \code{nume="valoare"}.

\begin{lstlisting}[style=html, caption={Tipuri de tag-uri}]
<!-- Tag container: are continut intre deschidere si inchidere -->
<p>Acesta este un paragraf.</p>

<!-- Tag void: nu are continut, nu se inchide -->
<img src="poza.jpg" alt="Descriere">

<!-- Tag cu atribute multiple -->
<a href="https://best-is.ro" target="_blank" rel="noopener">BEST</a>
\end{lstlisting}

Tag-urile void cuprind: \code{<img>}, \code{<br>}, \code{<hr>}, \code{<input>}, \code{<meta>}, \code{<link>}.

% =========================================================
\section{Metadate obligatorii in <head>}

Trei elemente sunt considerate obligatorii in orice pagina moderna:

\begin{enumerate}
  \item \code{<meta charset="UTF-8">} \textendash{} setul de caractere. UTF-8 acopera toate alfabetele moderne, inclusiv diacriticele romanesti.
  \item \code{<meta name="viewport"...>} \textendash{} controleaza redarea pe dispozitive mobile. Valoarea \code{width=device-width, initial-scale=1.0} instructeaza browserul sa pastreze latimea paginii egala cu cea a ecranului si sa nu zoom-eze initial.
  \item \code{<title>} \textendash{} apare in tab-ul browserului, in bookmark-uri, si ca titlu in rezultatele cautarii.
\end{enumerate}

\begin{atentie}
Conform datelor StatCounter (2024), peste 60\% din traficul web global provine de pe dispozitive mobile. Lipsa meta tag-ului \code{viewport} afecteaza direct experienta majoritatii utilizatorilor, generand o pagina nescalata, cu zoom necesar pentru lectura.
\end{atentie}

% =========================================================
\section{Tag-uri semantice vs <div> soup}

Tag-urile semantice descriu \emph{rolul} continutului, nu doar aspectul vizual. Browserul, motoarele de cautare si tehnologiile asistive (cititoare de ecran, navigatoare prin tastatura) pot interpreta automat structura paginii cand tag-urile semantice sunt folosite.

\begin{dano}[Asa nu \textendash{} ,,div soup'']
\begin{lstlisting}[style=html, numbers=none]
<div class="header">
  <div class="logo">Site</div>
  <div class="menu">
    <div class="link"><a href="/">Acasa</a></div>
  </div>
</div>
<div class="content">
  <div class="hero">
    <div class="title">Salut</div>
  </div>
</div>
<div class="footer">2026</div>
\end{lstlisting}
\textbf{Probleme.} Browserul si cititoarele de ecran nu stiu ce reprezinta fiecare \code{<div>}. Navigatia rapida prin landmarks este imposibila. SEO penalizat pentru lipsa de structura.
\end{dano}

\begin{dado}[Asa da \textendash{} semantic]
\begin{lstlisting}[style=html, numbers=none]
<header>
  <a class="logo" href="/">Site</a>
  <nav>
    <a href="/">Acasa</a>
  </nav>
</header>
<main>
  <section class="hero">
    <h1>Salut</h1>
  </section>
</main>
<footer>2026</footer>
\end{lstlisting}
\textbf{Beneficii.} Cititoarele de ecran ofera comenzi rapide pentru sarit la \code{<header>}, \code{<nav>}, \code{<main>}, \code{<footer>}. Google identifica zonele principale fara analiza euristica. Codul este mai usor de citit dupa 6 luni.
\end{dado}

\subsection{Repertoriul tag-urilor semantice}

\begin{itemize}[itemsep=2pt]
  \item \code{<header>} \textendash{} antetul unei pagini sau al unei sectiuni.
  \item \code{<nav>} \textendash{} bara de navigatie principala.
  \item \code{<main>} \textendash{} continutul unic al paginii (un singur \code{<main>} per pagina).
  \item \code{<section>} \textendash{} grup tematic de continut, de obicei cu un \code{<h2>}.
  \item \code{<article>} \textendash{} unitate de continut auto-suficienta (post de blog, comentariu, card produs).
  \item \code{<aside>} \textendash{} continut tangential (sidebar, note de subsol vizuale).
  \item \code{<footer>} \textendash{} subsolul unei pagini sau al unei sectiuni.
  \item \code{<figure>} + \code{<figcaption>} \textendash{} imagine cu legenda.
\end{itemize}

% =========================================================
\section{Ierarhia titlurilor}

Specificatia HTML5 defineste sase niveluri de titluri, de la \code{<h1>} la \code{<h6>}. Cititoarele de ecran genereaza un \emph{outline} al paginii pe baza acestei ierarhii; un utilizator nevazator parcurge documentul saltand intre titluri exact ca un cititor obisnuit cand parcurge un cuprins.

\begin{itemize}
  \item Un singur \code{<h1>} per pagina (conven\textcommabelow{t}ia industriala pentru accesibilitate \textcommabelow{s}i SEO; HTML Living Standard permite tehnic mai multe, dar Lighthouse \textcommabelow{s}i cititoarele de ecran a\textcommabelow{s}teapta unul singur), descriind tema centrala.
  \item Subsectiunile principale folosesc \code{<h2>}; sub-subsectiunile, \code{<h3>}.
  \item Nivelurile nu se sar (\code{<h1>} urmat direct de \code{<h3>} este o eroare structurala).
  \item Marimea vizuala se controleaza prin CSS, nu prin alegerea unui nivel diferit.
\end{itemize}

\begin{dano}[Asa nu \textendash{} stilul ales prin nivel]
\begin{lstlisting}[style=html, numbers=none]
<h1>Titlul mare</h1>
<h4>Subtitlu, am ales h4 ca arata mai mic</h4>
<h2>Sectiune</h2>
\end{lstlisting}
\end{dano}

\begin{dado}[Asa da \textendash{} stilul ales prin CSS]
\begin{lstlisting}[style=html, numbers=none]
<h1>Titlul mare</h1>
<h2 class="subtitle">Subtitlu</h2>
<h2>Sectiune</h2>
\end{lstlisting}
\code{.subtitle \{ font-size: 1rem; font-weight: normal; \}}
\end{dado}

% =========================================================
\section{Linkuri si navigare}

Elementul \code{<a>} (anchor) creeaza un hyperlink. Atributul \code{href} accepta mai multe forme:

\begin{lstlisting}[style=html, caption={Tipuri de URL-uri}]
<!-- URL absolut -->
<a href="https://best-is.ro">BEST Iasi</a>

<!-- URL relativ -->
<a href="contact.html">Contact</a>
<a href="../images/logo.png">Logo</a>

<!-- Ancora interna -->
<a href="#contact">Sari la contact</a>

<!-- Adresa email / numar de telefon -->
<a href="mailto:hello@best-is.ro">Email</a>
<a href="tel:+40123456789">Telefon</a>
\end{lstlisting}

\subsection{Linkuri externe \textendash{} target si rel}

Deschiderea unui link extern intr-un tab nou se face cu \code{target="\_blank"}. Fara protectie, pagina destinatie poate manipula tab-ul sursa prin \code{window.opener} (tabnabbing $\to$ phishing). Adauga mereu \code{rel="noopener noreferrer"}:

\begin{lstlisting}[style=html, numbers=none]
<a href="https://extern.ro" target="_blank" rel="noopener noreferrer">Extern</a>
\end{lstlisting}

\code{noopener} blocheaza accesul la \code{window.opener}; \code{noreferrer} ascunde URL-ul sursa.

% =========================================================
\section{Formulare accesibile}

Fiecare element \code{<input>} trebuie asociat cu un \code{<label>}. Asocierea se realizeaza prin perechea atribut-valoare \code{for}/\code{id}: valoarea atributului \code{for} de pe \code{<label>} este identica cu valoarea atributului \code{id} de pe \code{<input>}.

\begin{dano}[Asa nu \textendash{} placeholder ca substitut pentru label]
\begin{lstlisting}[style=html, numbers=none]
<input type="email" placeholder="Email">
\end{lstlisting}
\textbf{Probleme.} (1) Cititoarele de ecran nu anunta placeholder-ul ca eticheta. (2) Placeholder-ul dispare la tastare, lasand utilizatorul fara context. (3) Contrast scazut pe textul de placeholder.
\end{dano}

\begin{dado}[Asa da]
\begin{lstlisting}[style=html, numbers=none]
<label for="email">Email</label>
<input type="email" id="email" name="email" required>
\end{lstlisting}
\textbf{Beneficii.} Cititoarele anunta corect; clic-ul pe text focuseaza inputul; eticheta ramane vizibila.
\end{dado}

\subsection{Tipuri de input frecvente}

Atributul \code{type} controleaza interfata oferita de browser si validarea. Pe mobil, tipul determina tastatura afisata.

\begin{itemize}[itemsep=2pt]
  \item \code{type="text"} \textendash{} text generic.
  \item \code{type="email"} \textendash{} validare format email; tastatura cu \code{@} pe mobil.
  \item \code{type="tel"} \textendash{} tastatura numerica pe mobil; nu valideaza formatul (variaza pe tari).
  \item \code{type="number"} \textendash{} accepta doar cifre; arrows pe desktop.
  \item \code{type="date"} \textendash{} date picker nativ.
  \item \code{type="password"} \textendash{} ascunde caracterele.
  \item \code{type="url"} \textendash{} validare format URL.
  \item \code{type="search"} \textendash{} buton clear in unele browsere.
\end{itemize}

\begin{ponturi}
Browserul valideaza nativ inputurile: campurile cu atributul \code{required} blocheaza submit-ul daca sunt goale; \code{type="email"} respinge valori care nu contin \code{@}. Aceasta validare \textbf{nu inlocuieste} validarea pe server, dar imbunatateste experienta utilizatorului inainte de a trimite cererea.
\end{ponturi}

\subsection{Element HTML \textendash{} rol UI generic}

Daca ai mai folosit o interfata grafica (QtCreator, Figma, Visual Studio, Android Studio), recunosti aceste pattern-uri sub alte nume. Tabelul de mai jos pune in oglinda terminologia:

\begin{tcolorbox}[
  enhanced, breakable, sharp corners, arc=0pt,
  colback=accentSoft, colframe=partcolor, boxrule=0pt, leftrule=3pt,
  left=12pt, right=12pt, top=8pt, bottom=8pt,
  fontupper=\small,
]
\renewcommand{\arraystretch}{1.35}
\begin{tabular}{@{}p{0.42\linewidth}@{\hspace{8pt}}p{0.52\linewidth}@{}}
{\sffamily\bfseries\color{partcolor} Element HTML} & {\sffamily\bfseries\color{partcolor} Rolul UI generic} \\
\code{<label>}                       & text static (eticheta unui camp) \\
\code{<input type="text">}           & lineedit -- input pentru text scurt \\
\code{<input type="email">}          & lineedit cu validare email \\
\code{<input type="password">}       & lineedit ascuns (parola) \\
\code{<input type="number">}         & spinbox -- numar cu sageti \\
\code{<input type="date">}           & date picker \\
\code{<input type="range">}          & slider -- valoare pe interval \\
\code{<textarea>}                    & text multi-linie \\
\code{<select>} + \code{<option>}    & combobox -- o alegere dintr-o lista \\
\code{<input type="radio">}          & radio button -- o singura alegere dintr-un grup \\
\code{<input type="checkbox">}       & checkbox -- alegeri multiple (sau pornit/oprit) \\
\code{<input type="file">}           & file dialog -- selectie fisier \\
\code{<input type="color">}          & color picker \\
\code{<button>}                      & buton de actiune \\
\code{<a href="...">}                & link -- navigare catre URL \\
\code{<img>}                         & imagine statica \\
\code{<video>} / \code{<audio>}      & player media nativ \\
\code{<details>} + \code{<summary>}  & accordeon -- arata/ascunde \\
\code{<dialog>}                      & modal -- fereastra suprapusa \\
\code{<progress>}                    & bara de progres \\
\end{tabular}
\end{tcolorbox}

Regula: cand alegi un element, te intrebi mai intai ,,ce rol UI joaca?'' \textendash{} nu ,,cum ar arata cu CSS''. Tag-ul corect rezolva accesibilitatea, validarea \textcommabelow{s}i comportamentul nativ \emph{gratis}.

% =========================================================
\section{Atribute esentiale}

\begin{itemize}
  \item \code{href} pe \code{<a>} \textendash{} URL-ul tinta.
  \item \code{src} pe \code{<img>} \textendash{} calea catre imagine.
  \item \code{alt} pe \code{<img>} \textendash{} text descriptiv. Obligatoriu conform WCAG 2.1 (criteriul 1.1.1 \emph{Non-text Content}).
  \item \code{type} pe \code{<input>} \textendash{} determina validarea si interfata.
  \item \code{required} pe \code{<input>} \textendash{} activeaza validarea nativa.
  \item \code{lang} pe \code{<html>} \textendash{} limba documentului.
  \item \code{title} pe orice element \textendash{} tooltip la hover (uz limitat, nu inlocuieste \code{aria-label}).
\end{itemize}

\subsection{Calitatea textului \texttt{alt}}

Textul \code{alt} descrie ce \emph{vede} imaginea, in contextul paginii. Scopul: utilizatorul fara acces vizual primeste informatia echivalenta.

\begin{dano}[Asa nu \textendash{} alt nesemnificativ sau spam SEO]
\begin{lstlisting}[style=html, numbers=none]
<img src="echipa.jpg" alt="">
<img src="echipa.jpg" alt="poza">
<img src="echipa.jpg" alt="best best best iasi training web html css">
\end{lstlisting}
\end{dano}

\begin{dado}[Asa da \textendash{} alt descriptiv si specific]
\begin{lstlisting}[style=html, numbers=none]
<img src="echipa.jpg" alt="Echipa BEST Iasi in fata sediului UAIC, primavara 2026">
\end{lstlisting}
\textbf{Exceptie:} pentru imagini pur decorative (fundal, separator, ornament), \code{alt=""} este corect \textendash{} cititorul de ecran le ignora.
\end{dado}

% =========================================================
\section{Tag-uri suplimentare \textendash{} men\textcommabelow{t}iuni rapide}

Pentru proiecte mai complexe, HTML5 ofera elemente specializate care nu intra in scopul fazei curente, dar merita cunoscute la nivel de existen\textcommabelow{t}a:

\begin{itemize}[itemsep=2pt]
  \item \code{<table>} cu \code{<thead>}, \code{<tbody>}, \code{<tfoot>}, \code{<caption>} \textendash{} pentru date tabulare reale (nu pentru layout). Atributele \code{scope}, \code{headers} asigura accesibilitatea.
  \item \code{<video>} si \code{<audio>} \textendash{} mediafile native, fara nevoie de plugin-uri (Flash, Silverlight au murit). Atribute: \code{controls}, \code{autoplay}, \code{muted}, \code{loop}, \code{poster} (pentru video).
  \item \code{<canvas>} \textendash{} suprafa\textcommabelow{t}a 2D pentru desenare programatica cu JavaScript. Folosit pentru jocuri, vizualizari, anima\textcommabelow{t}ii complexe.
  \item \code{<details>} si \code{<summary>} \textendash{} accordeon nativ, fara JavaScript.
  \item \code{<dialog>} \textendash{} modal nativ, suportat in toate browserele moderne.
  \item \code{<picture>} \textendash{} container pentru \code{<img>} cu surse alternative pe formate sau dimensiuni.
\end{itemize}

Documenta\textcommabelow{t}ia exhaustiva: MDN \textendash{} HTML elements reference (link in resursele de la finalul capitolului).

% =========================================================
\section{Comentarii si caractere speciale}

Comentariile HTML nu apar in pagina dar raman vizibile in \emph{View Source}. Sintaxa: \code{<!-- text -->}.

Caracterele rezervate (\code{<}, \code{>}, \code{\&}) trebuie scrise prin entitati pentru a nu fi interpretate ca markup:

\begin{itemize}[itemsep=1pt]
  \item \code{\&lt;} \textendash{} \code{<}
  \item \code{\&gt;} \textendash{} \code{>}
  \item \code{\&amp;} \textendash{} \code{\&}
  \item \code{\&copy;} \textendash{} \code{(c)}
  \item \code{\&nbsp;} \textendash{} spatiu nedespartibil
\end{itemize}

% =========================================================
\begin{bonusbox}
\textbf{Acest box e optional.} Daca esti la inceput, treci direct la \emph{Verificare}. Pattern-urile de mai jos preg\u atesc terenul pentru Lighthouse \textcommabelow{s}i accesibilitate avansata.

\medskip

\textbf{\faImage\ \texttt{<picture>} \textendash{} imagini responsive cu fallback}

\code{<img>} obi\textcommabelow{s}nuit serve\textcommabelow{s}te o singura sursa. \code{<picture>} alege intre mai multe surse: format modern (AVIF $\to$ WebP $\to$ JPEG fallback) sau dimensiune (mobil vs desktop):

\begin{lstlisting}[style=html, numbers=none]
<picture>
  <source srcset="hero.avif" type="image/avif">
  <source srcset="hero.webp" type="image/webp">
  <img src="hero.jpg" alt="Sediul BEST Iasi" width="800" height="450">
</picture>
\end{lstlisting}

Browserul incarca \textbf{prima sursa pe care o sus\textcommabelow{t}ine}. Chrome modern $\to$ AVIF (cel mai mic), Safari vechi $\to$ JPEG. Zero JavaScript, zero detec\textcommabelow{t}ie de feature.

\medskip

\textbf{\faBolt\ \texttt{loading="lazy"} \textcommabelow{s}i \texttt{decoding="async"}}

Doua atribute care imbunata\textcommabelow{t}esc semnificativ Lighthouse fara cost:

\begin{lstlisting}[style=html, numbers=none]
<img src="profile.jpg" alt="Voluntar BEST"
     loading="lazy"           <!-- amana download pana cand intra in viewport -->
     decoding="async"         <!-- decodifica pe alt thread, nu blocheaza UI -->
     width="200" height="200">
\end{lstlisting}

\textbf{Excep\textcommabelow{t}ie:} imaginea hero (prima vizibila) \textbf{nu} primeste \code{loading="lazy"} \textendash{} ar intarzia LCP. Pentru ea: \code{fetchpriority="high"}.

\medskip

\textbf{\faUniversalAccess\ ARIA \textendash{} c\u and semantica nu e suficienta}

ARIA (\emph{Accessible Rich Internet Applications}) adauga semnificatie pentru cititoarele de ecran cand HTML-ul semantic nu acopera contextul. \textbf{Regula de aur:} prima alegere e mereu un tag semantic. ARIA il completeaza.

Cele 3 atribute pe care le folose\textcommabelow{s}ti des:

\begin{itemize}[itemsep=2pt]
  \item \code{aria-label="..."} \textendash{} eticheta accesibila cand textul vizibil lipseste sau e ambiguu. Exemplu clasic: butonul hamburger.
\begin{lstlisting}[style=html, numbers=none]
<button class="menu-toggle" aria-label="Deschide meniul">
  <svg>...</svg>
</button>
\end{lstlisting}

  \item \code{aria-current="page"} \textendash{} indica link-ul activ in navigare:
\begin{lstlisting}[style=html, numbers=none]
<nav>
  <a href="/" aria-current="page">Acasa</a>
  <a href="/projects">Proiecte</a>
</nav>
\end{lstlisting}

  \item \code{role="..."} \textendash{} \textbf{rar necesar}. Folose\textcommabelow{s}te-l doar cand transformi un \code{<div>} in ceva semantic (ce \textbf{nu ar trebui sa faci}). Excep\textcommabelow{t}ie utila: \code{role="alert"} pentru notificari live.
\end{itemize}

\begin{atentie}
NU adauga \code{role="button"} pe un \code{<button>}, sau \code{role="navigation"} pe un \code{<nav>}. Tag-ul semantic are deja rolul corect. ARIA duplicat = bruiaj pentru cititoarele de ecran.
\end{atentie}

\medskip

\textbf{\faRoute\ Unde mergi de aici}

\begin{itemize}[itemsep=1pt]
  \item \emph{MDN \texttt{<picture>} reference:} \url{https://developer.mozilla.org/en-US/docs/Web/HTML/Element/picture}
  \item \emph{web.dev Lazy load images:} \url{https://web.dev/articles/browser-level-image-lazy-loading}
  \item \emph{WAI-ARIA Authoring Practices:} \url{https://www.w3.org/WAI/ARIA/apg/}
\end{itemize}
\end{bonusbox}

% =========================================================
\section*{Verificare cuno\textcommabelow{s}tin\textcommabelow{t}e}

\begin{selfcheck}
\begin{enumerate}[itemsep=2pt]
  \item Care este diferen\textcommabelow{t}a func\textcommabelow{t}ionala intre \code{<div class="header">} \textcommabelow{s}i \code{<header>}?
  \item De ce este nevoie de \code{<meta name="viewport">} pentru afi\textcommabelow{s}area pe mobil?
  \item \emph{Spot the issue:} \code{<input type="email" placeholder="Email">}.
  \item Cate elemente \code{<h1>} recomanda conven\textcommabelow{t}ia industriala (accesibilitate + SEO) sa existe intr-o pagina?
  \item Care sunt cele doua atribute necesare cand un link extern se deschide intr-un tab nou?
\end{enumerate}
\end{selfcheck}

% =========================================================
\section*{Recapitulare}

\begin{recap}[Faza 1 -- HTML]
\begin{itemize}[itemsep=1pt]
  \item Anatomia: \code{<!DOCTYPE html>} $\rightarrow$ \code{<html lang>} $\rightarrow$ \code{<head>} (metadate) $\rightarrow$ \code{<body>} (continut).
  \item Trei meta tags obligatorii: \code{charset}, \code{viewport}, \code{<title>}.
  \item Tag-uri semantice in locul \code{<div>}: \code{<header>}, \code{<main>}, \code{<section>}, \code{<footer>}, \code{<nav>}, \code{<article>}, \code{<aside>}.
  \item Un singur \code{<h1>} per pagina; ierarhia titlurilor nu admite salturi de nivel.
  \item Forms accesibile: \code{<label for="X">} corespunde cu \code{<input id="X">}.
  \item Atribute obligatorii pentru accesibilitate: \code{alt} pe \code{<img>}, \code{lang} pe \code{<html>}.
  \item Linkuri externe in tab nou: combinatia \code{target="\_blank"} + \code{rel="noopener noreferrer"}.
  \item \faRocket\ Bonus: \code{<picture>} pentru fallback AVIF/WebP/JPEG, \code{loading="lazy"} + \code{decoding="async"} pe imagini, ARIA (\code{aria-label}, \code{aria-current}).
\end{itemize}
\end{recap}

\partdivider

% =========================================================
\begin{joc}
\textbf{Platforma:} CodePen (\url{codepen.io}). \quad \textit{URL-ul exact este distribuit de trainer in chat.}

\medskip

\textbf{Activitatea.} Documentul de pornire contine o structura non-semantica (\code{<div>} folosit excesiv), atribute \code{alt} lipsa pe imagini si elemente \code{<input>} fara \code{<label>} asociat. Sarcina este de a transforma documentul intr-o varianta semantica si conformanta cu standardele de accesibilitate.

\medskip

\textbf{Criterii de validare:}
\begin{itemize}[itemsep=1pt]
  \item Folosirea tag-urilor \code{<header>}, \code{<main>}, \code{<section>}, \code{<footer>} in locul \code{<div>}.
  \item Prezenta atributului \code{alt} pe toate elementele \code{<img>}.
  \item Asocierea \code{<label for>} \textendash{} \code{<input id>} pentru fiecare camp de formular.
  \item Existenta unui singur element \code{<h1>} per pagina.
  \item Linkuri externe au \code{target="\_blank"} + \code{rel="noopener noreferrer"}.
\end{itemize}

\medskip

\textbf{Durata estimata:} 5 minute.
\end{joc}

% =========================================================
\section*{\faGraduationCap\ \ Resurse pentru aprofundare}

\begin{itemize}[itemsep=2pt]
  \item Mozilla Developer Network. \emph{HTML elements reference}. \\\url{https://developer.mozilla.org/en-US/docs/Web/HTML/Element}
  \item The Odin Project. \emph{HTML Foundations}. \\\url{https://www.theodinproject.com/paths/foundations/courses/foundations\#html-foundations}
  \item Google. \emph{Learn HTML}. \url{https://web.dev/learn/html}
  \item HTML Reference. \emph{Vizual cheatsheet de tag-uri}. \url{https://htmlreference.io}
  \item Google. \emph{Learn Forms}. \url{https://web.dev/learn/forms}
  \item W3C. \emph{Web Accessibility Initiative \textendash{} WCAG}. \url{https://www.w3.org/WAI/standards-guidelines/wcag/}
\end{itemize}

\bigskip

\bridge{Capitolul 2 (CSS) imbraca structura HTML cu stil vizual. Aspectul, layout-ul \textcommabelow{s}i responsivitatea se construiesc plecand de la pagina semantica generata aici.}
